%!TEX root = ../main.tex

\chapter{Prompt finale per Prism}

\section{Prompt definitivo}

\begin{quote}\small
Genera una relazione finale di tirocinio universitario in italiano, con tono accademico, chiaro, ordinato e professionale. La relazione deve avere un corpo centrale di circa dieci pagine, senza contare copertina, indice, bibliografia e appendici.

Il progetto si chiama \builder{} ed è un'applicazione web per la progettazione, visualizzazione, traduzione ed esportazione di schemi Entity-Relationship e schemi logici. La relazione deve essere una vera relazione di tirocinio, non una documentazione tecnica e non un manuale utente.

Usa tutti i materiali preparati nelle Fasi 1--6: raccolta informazioni, outline, bozza completa, integrazione del ruolo della tutor, revisione stilistica e versione finale consegnabile. Mantieni coerenza terminologica e riduci le ripetizioni.

La tutor del tirocinio è stata la Prof.ssa Alessandra Lumini. Valorizza il suo ruolo di supervisione, spiegando che le sue indicazioni hanno contribuito a orientare il lavoro verso chiarezza concettuale, correttezza degli schemi, leggibilità grafica, usabilità, valore didattico, manutenibilità e revisione progressiva. Non inventare citazioni dirette e non attribuire alla tutor frasi testuali non documentate.

La relazione deve includere: introduzione; obiettivi del tirocinio; descrizione del progetto \builder; analisi dei requisiti; tecnologie utilizzate; architettura e progettazione; funzionalità sviluppate e consolidate; test, verifica e miglioramenti; supervisione della tutor e competenze acquisite; conclusioni e sviluppi futuri.

Usa le informazioni tecniche verificate: React, TypeScript, Vite, CSS, canvas SVG, npm, \texttt{tsx}, Playwright, Git/GitHub, GitHub Actions, GitHub Pages, \texttt{lucide-react}, export PNG/SVG/JPEG, formato \texttt{.ersp}, sincronizzazione ERS, traduzione verso schema logico, generazione SQL, reverse engineering SQL, localizzazione in italiano, inglese e albanese, test unitari, test di integrazione e test end-to-end.

Non inventare dati mancanti. Dove mancano informazioni amministrative, usa placeholder chiari, ad esempio \placeholder{Periodo del tirocinio da completare}, \placeholder{Ore/CFU da completare}, \placeholder{Corso di laurea da completare}, \placeholder{Matricola da completare}, \placeholder{Anno accademico da completare}.

Mantieni equilibrio tra parte tecnica e parte formativa. Inserisci tabelle dove utili: obiettivi, tecnologie, test e impatto della supervisione della tutor. Suggerisci punti in cui inserire screenshot dell'interfaccia, dello schema ER, della vista logica e dell'export.

Prepara un testo adatto a PDF, Word o LaTeX, con paragrafi ordinati, titoli chiari, limiti realistici e sviluppi futuri credibili.
\end{quote}
